\documentclass[twocolumn]{aastex631}
\usepackage{booktabs}
%% Float placement: let full-width figure*/deluxetable* sit at page tops with
%% text below (minimal white space), and only form a dedicated float page when
%% a float genuinely fills most of it.
\setcounter{topnumber}{3}
\setcounter{dbltopnumber}{3}
\setcounter{totalnumber}{5}
\renewcommand{\topfraction}{0.9}
\renewcommand{\dbltopfraction}{0.9}
\renewcommand{\textfraction}{0.06}
%% keep floatpagefraction high so a wide figure is placed at a page top WITH
%% text below it, instead of getting a near-empty dedicated float page.
\renewcommand{\floatpagefraction}{0.9}
\renewcommand{\dblfloatpagefraction}{0.9}
\shorttitle{An engine-regulated covering law for LRDs}
\shortauthors{Pandev}

\begin{document}

\title{An Engine-Regulated Covering Law for Little Red Dots: One Equilibrium
Relation for Their Spectral Sequence, Demography, and Broad-Line Widths}

\author[0009-0005-8153-071X]{Stilian Pandev}
\affiliation{Independent researcher}
\email{stilianpandev@gmail.com}

\correspondingauthor{Stilian Pandev}

\begin{abstract}
Little red dots (LRDs) present three puzzles usually treated separately: an
ordered spectral-subtype sequence, a rise--peak--decline number density,
and broad Balmer lines in compact, X-ray-weak sources. Within a
semi-analytic lifecycle model with a dense-plus-diffuse soft-EUV emission
layer, we show that all three respond to a single equilibrium relation: the
covering odds of the nuclear reprocessor equal the square of the intrinsic
engine-to-host continuum contrast, $C/(1-C)=x^{2}$ with $x\equiv
L_{\nu,\rm engine}(2500\,\text{\AA})/L_{\nu,\star}(2500\,\text{\AA})$ and
unit normalization. Confronted directly with the eight stacked
[O~III]/H$\beta$ and [O~III]/[O~II] ratios of the four continuum subtypes,
the raw ionizing-photon partition under this law---supplemented by one
physically motivated geometric parameter, the stellar ionizing coupling
$\epsilon\simeq0.1$, consistent with canonical Lyman-continuum escape
fractions---matches the data at least as well as an imposed
photon-partition closure ($\chi^{2}=10.4$ versus $12.9$ on the identical
metric; seed range 5--16), while the demographic fit on the reliable
$z\le7.5$ bins improves from $\chi^{2}=6.9$ to $1.2$ and the RUBIES
broad-width distribution distance improves from 0.41 to 0.35 with the
widths reinterpreted as substantially electron-scattering broadened
($\tau_{\rm e}\simeq2.3$). The exponent is not fitted but constructed and
then confronted: of four candidate mechanisms, roll-over (sequential-stage)
and steepened (Poisson double-blocking) laws are rejected by the data,
cooling-fractionation is excluded both by its saturated timescales and by
the demonstrated metallicity-blindness of the law
($\Delta\chi^{2}>120$ against any $Z$-modulation), leaving a unique
survivor among the candidates considered---momentum-limited envelope
building against drag-limited clump removal---whose precondition we
verify holds for 100\% of the selected population, and whose structure
identifies $\gamma=2$ with the number of soft modes of the doubly forced
clump system. The model predicts a subtype
``march'' (selected red fraction $0.04\pm0.05$ at $z\simeq9.5$ rising to
$0.49\pm0.09$ at $z\simeq4.5$) and blue LRDs that are more X-ray
transmitting and more variable than red LRDs at matched luminosity. We
also show the reported $z\simeq9.5$ abundance cannot currently constrain
formation: transparent variations of one color-selection gate change the
predicted density by an order of magnitude while moving the $z\le7.5$ bins
by well under their errors.
\end{abstract}

\keywords{Active galactic nuclei (16) --- Supermassive black holes (1663)
--- High-redshift galaxies (734) --- Galaxy evolution (594)}

\section{Introduction}\label{sec:intro}
JWST has revealed an abundant population of compact, red, frequently
broad-line sources at $z\gtrsim4$, the little red dots
\citep{Matthee2024,Greene2024,Kocevski2024,Akins2024,Kocevski2025}. Three
empirical regularities frame the present work. First, stacked spectroscopy
resolves an ordered continuum-subtype sequence, from extreme red (xLRD)
through blue (bLRD) classes, with common line species but strongly varying
line-to-continuum structure \citep{PerezGonzalez2026}, a diversity also
mapped to black-hole luminosity and host contrast in a stationary
framework \citep{Barro2025}. Second,
broad-selection number densities rise, peak near $z\sim5$--6, and decline
toward both lower and higher redshift \citep{Rinaldi2026,Ma2026}. Third,
the population combines broad Balmer lines in very compact regions with
pronounced X-ray weakness
\citep{Ananna2024,Yue2024,Maiolino2024,Sneppen2026} and generally low
photometric variability \citep[][the last for a five-object
sample]{Tee2024,Secunda2025,KokuboHarikane2024},
plausibly indicating dense, Compton-thick circumnuclear gas. A wave of
recent models interprets LRDs as black holes embedded in dense gas
envelopes with quasi-stellar photospheres and scattering-broadened Balmer
lines \citep{Naidu2025,Rusakov2025,deGraaff2025,Sneppen2026b}, a programme
since extended to unified ``black-hole-star'' families and explicit
nuclear-minus-host decompositions \citep{deGraaff2025b,Sun2026}
and to quasi-star evolutionary tracks and late-stage quasi-star models
\citep{Santarelli2025,Begelman2026,Kido2025};
super-Eddington accretion states offer an alternative account of the
X-ray weakness \citep{PacucciNarayan2024,Inayoshi2024}, though deep
Chandra limits challenge some variants \citep[e.g.,][]{Sacchi2025}; see
the current X-ray census of \citet{Tortosa2026}.
Published clustering measurements additionally bound the host halos
\citep{Pizzati2025,Carranza2025,Wang2026}, a constraint the lifecycle's
halo occupation must ultimately face. Our model shares the
dense-envelope architecture and the low-spin entry channel proposed for
LRD demographics by \citet{PacucciLoeb2025}; its distinguishing content
is a population-level equilibrium law connecting the envelope to the
host, and the joint confrontation of that law with the subtype sequence,
the demography, and the width distribution simultaneously.

In a companion demographic analysis (Pandev 2026a, in preparation) we
concluded that binned number densities alone cannot select a formation
mechanism. Here we take the opposite approach: a forward semi-analytic
lifecycle model whose observation layer is confronted simultaneously with
the stacked spectral sequence, the number densities, public broad-line
widths \citep{Hviding2025}, and population-level hardness constraints
\citep{Sok2026}. The central result is that these independent
confrontations select, and are jointly improved by, one simple relation
between the covering of the nuclear reprocessor and the visibility of the
nuclear engine---and that the relation's form can be constructed from an
explicit elimination program rather than fitted.

Section~\ref{sec:model} summarizes the model.
Section~\ref{sec:closure} develops the covering law and its derivation
status. Section~\ref{sec:results} presents the confrontations,
Section~\ref{sec:predictions} the predictions, Section~\ref{sec:highz} the
status of the high-redshift bin, and Section~\ref{sec:limits} the
limitations.

\section{Model}\label{sec:model}

\subsection{Lifecycle}
The population model evolves $1.2\times10^{4}$ halo histories, importance
reweighted to the \citet{Tinker2008} mass function, with mean growth from
\citet{Fakhouri2010} and halo bias from \citet{Tinker2010}. Each halo
carries nuclear gas, black-hole, and burst-reservoir masses through the
states quiescent $\rightarrow$ compacting $\rightarrow$ embedded
$\rightarrow$ clearing $\rightarrow$ post-LRD, with recurrence. Entry
depends on spin \citep{Bullock2001}, specific growth, and nuclear Thomson
depth \citep[low-angular-momentum compaction;][]{Mo1998}; embedded objects
exit through competing gas-exhaustion, AGN-feedback, and host/size-growth
hazards. Population survival is emergent: no survival kernel is imposed. A
secondary gas-rich burst entry channel operates in rapidly assembling
high-redshift hosts; it transfers a finite baryon reservoir and is
mass-conserving. One global abundance normalization is calibrated at
$4.5<z<6.5$ and is refit for each model configuration; comparative
statistics below therefore compare shapes, not absolute abundances.

\subsection{Observation layer}
The porous-envelope layer assigns each active source a two-scale geometry:
a cool reprocessor of outer scale 1360 AU with covering fraction $C$,
clump-population UV/X-ray transfer, and a compact line-emitting scale
calibrated in Section~\ref{sec:widths}. The forward continuum is the sum
of an attenuated stellar host, a leaked accretion component, and the
reprocessor photosphere; sources are classified into the four continuum
subtypes by the published $L_{5100}/L_{2500}$ boundaries (1.8, 3.1, 6.3
from bLRD through xLRD), and the ``V-shape color'' below denotes the
model's rest-frame optical-redness proxy used by its selection function
(the empirical V-shaped turnover at the Balmer limit is characterized by
\citealt{Setton2024}).
The emission layer is a dense-plus-diffuse two-zone photoionization model
built on thermal-balance Cloudy anchors \citep[C23 release; a locally
verified C25 build is used;][]{Chatzikos2023}: a dense, dust-free skin
producing H/He lines with collisionally suppressed forbidden lines, and a
diffuse skin producing the optical metal lines, mixed at the level of
photon yields by the dense fraction of intercepted ionizing photons,
$f_{\rm dense}$. Public stacks reject a standard AGN ionizing continuum
through He~II: the data require a soft basis that preserves 35--54 eV
photons while suppressing flux above the He$^{+}$ edge; a ten-screen
warm-absorber scan shows a generic ionized filter cannot achieve this, so
a soft basis (labelled by a 50 kK spectral shape, not a literal
temperature) is adopted for all sources
\citep{PerezGonzalez2026,Sok2026}.

\section{The covering law and its derivation status}\label{sec:closure}

\subsection{Decomposition of the imposed closure}\label{sec:decomp}
Inverting the two-zone model through the four stacked subtype spectra
yields an ordered sequence in the dense fraction of processed ionizing
photons: 0.981, 0.901, 0.738, and 0.086 (xLRD$\rightarrow$bLRD; intervals
in Table~\ref{tab:ratios}). We emphasize that these are model-conditional
inversions---obtained through the same two-zone emission layer used
forward---not direct measurements; the directly observed quantities are
the stacked line ratios confronted in Section~\ref{sec:spectra}. Defining
the nuclear-to-host continuum dominance
$D\equiv L_{\nu,\rm thermal}(5100)/L_{\nu,\rm host}(5100)$, the inverted
sequence is reproduced by the quadratic closure
$f_{\rm dense}=D^{2}/(1+D^{2})$, with a free exponent returning
$2.02$ (conditional on a declared $0.05$ systematic floor discussed in
Section~\ref{sec:spectra}; $1.89$ with measurement errors alone).

To determine what supplies the two powers, we constructed every physical
photon partition available upstream of the closure---envelope-intercepted
engine photons against diffuse illumination by stellar, escaped-nuclear,
and dust-quenched supplies---and measured the emergent source-level
exponent of $\log(Q_{\rm dense}/Q_{\rm diffuse})$ on $\log D$
(Figure~\ref{fig:exponent}). Restricted to the physically populated range
$|\log_{10}D|<1.5$ (the extended panels also contain numerically capped
sources at extreme $D$, produced by a floor on the host term; they are
excluded from all fits and flagged in the caption), the fiducial
candidates give $\eta=0.65\pm0.10$ and $0.66\pm0.11$; the legacy
Thomson-depth covering contributes $0.09$ powers (slope of covering odds
on $\log D$). Ordinary flux competition therefore supplies roughly one of
the two required powers, and the second must come from co-regulation of
the interception itself. We record for completeness that our pre-declared
expectation---that installing a covering law would raise the emergent
source-level exponent to $\simeq2$---was not borne out (it rises only to
$0.65$--$0.89$); the resolution is that the stacked observables constrain
subtype-level medians, not the source-level slope, and the within-subtype
scatter of Section~\ref{sec:predictions} is the observable that tests the
difference.

\begin{figure*}[tp]
\centering
\includegraphics[width=0.98\textwidth]{interception_exponent.pdf}
\caption{Emergent source-level exponent of four physical photon partitions
against the continuum dominance $D$ at $4.5\le z\le6.5$. Fits (red, with
standard errors) are restricted to the physically populated range
$|\log_{10}D|<1.5$ (unshaded); the shaded regions contain numerically
capped sources (floors on the host term and yield caps) excluded from all
fits. Ordinary flux competition supplies roughly one power of $D$; the
second must come from the covering law.}
\label{fig:exponent}
\end{figure*}

\subsection{The covering law}\label{sec:law}
We therefore replace the Thomson-depth covering with an engine-regulated
law. With
$x\equiv L_{\nu,\rm engine}(2500\,\text{\AA})/
L_{\nu,\star}(2500\,\text{\AA})$
the \emph{intrinsic} (pre-attenuation, pre-covering) contrast, the
covering odds are
\begin{equation}
\frac{C}{1-C}=k\,x^{\gamma},
\label{eq:law}
\end{equation}
with a frozen per-halo lognormal scatter of $0.35$ dex (this value is an
assumption: a scan over $0.20$--$0.50$ dex changes the spectral statistic
by $<5$\%, so the scatter is not constrained by the present observables).
Both terms in $x$ are computed upstream of the observables inside the
simulator; we caution that this does not by itself make the calibration
non-circular, and we address circularity by construction: the law's
parameters are assessed against the directly observed line ratios
(Section~\ref{sec:spectra}) and the $z\le7.5$ demography, and the
model-conditional inversions of Section~\ref{sec:decomp} are used only as
diagnostics. A scan over $\gamma\in\{1.5,2,2.5,3\}$,
$k\in\{0.5,1,1.5,2\}$ (a descriptive sum of the spectral and demographic
statistics, not a likelihood) prefers $(\gamma,k)=(2,1)$; the preference
is stable across four random seeds and against removal of the contested
$z=9.5$ bin (Section~\ref{sec:highz}), and under the fiducial
line-ratio-space objective the nearest competitor---the $\gamma=3$,
$k=0.5$ point---sits at joint $\simeq32$ versus 11.5, a factor $\sim$3
worse; the tabulated $k=1$ profile is itself non-monotonic in $\gamma$
(Table~\ref{tab:elimination}), reflecting selection reshuffling between
laws. Covering odds
equal the squared contrast with 50\% covering at engine--host parity.

\begin{figure*}[tp]
\centering
\includegraphics[width=0.9\textwidth]{coregulated_covering_scan.pdf}
\caption{Joint scan of the covering law under the fiducial metric. Left:
joint objective (flux-stacked line-ratio $\chi^{2}$ plus $z\le7.5$
demographic $\chi^{2}$; logarithmic color scale) over the evaluated
$(\gamma,k)$ cells; the preferred point is the unit-normalization
quadratic. Right: spectral versus demographic components for all
evaluated points. The objective is a descriptive sum; no probabilistic
reading of $\Delta$ values is intended.}
\label{fig:scan}
\end{figure*}

\subsection{Equilibrium structure, anchors, and bounds}\label{sec:equil}
Three features of the law are derivable. (i) \emph{The odds form.} If
covering evolves by competition between a building rate per uncovered
solid angle, $B$, and a clearing rate per covered solid angle, $G$,
$\dot C=(1-C)B-CG$ equilibrates exactly at $C/(1-C)=B/G$. (ii) \emph{The
normalization.} Identifying $B$ with the engine momentum flux and $G$ with
host feedback, the ratio at a shared interface is radius-independent,
$B/G\sim(L_{e}/c)/(\dot p_{\rm SN}\,{\rm SFR})$ with
$\dot p_{\rm SN}\simeq3\times10^{33}$ dyn per $M_{\odot}\,{\rm yr^{-1}}$
\citep{KimOstriker2015}, and evaluates to $\simeq3$ at $x=1$: the anchor
argument therefore fixes $k$ to order unity within a factor $\sim3$, while
the scan discriminates at the factor-2 level; we present the anchor as a
consistency check at its actual precision, not as a derivation of $k=1$.
(iii) \emph{Fast relaxation.} A time-resolved implementation (exact
exponential-integrator update; covering initialized unbuilt at episode
entry) shows the stationary law is the fast-relaxation limit: relaxation
timescales of 30, 80, and 300 Myr degrade the spectral statistic
monotonically (defining the bound by where the degradation exceeds the
seed scatter gives $\tau_{\rm relax}\lesssim30$ Myr, the shortest value
resolvable at the 12 Myr integration step). The failure mode is fossil
covering: envelopes built while the engine was strong and retained as it
fades repopulate the blue subtype class with dense sources. The bound sits
six orders of magnitude above the $\sim$10 yr dynamical time at the
envelope scale, so fast equilibration is the expected regime.

\subsection{The exponent: a completed elimination program}
\label{sec:exponent}
The mass budget alone cannot supply the second power of $x$: converting
the engine's clump supply and the host's ablation flux to solid-angle
rates, the clump surface density cancels between build and clear, so any
mass-limited build against entrainment-limited clearing yields exactly one
power of $x$ regardless of confinement details. Four candidate mechanisms
for the second power were formulated and confronted
(Table~\ref{tab:elimination}):

\emph{Route B (sequential two-stage odds)}, $C/(1-C)=(kx)^{2}/(1+2kx)$
with the combinatorial structure forcing $k\simeq1$: rejected in both
forms under the fiducial metric (Table~\ref{tab:elimination}). The
parameter-free form fails catastrophically because its high-contrast
roll-over flattens the law across the populated range
($x\simeq0.5$--2); the profile minimum in $k$ ($k=3$, which would also
displace the per-stage parity to $x=1/3$) still loses to the pure
quadratic by a margin an order of magnitude beyond the seed scatter.

\emph{Route D (Poisson double-blocking)}, $C=1-e^{-kx}(1+kx)$ from
marginally translucent clumps requiring two-deep sightlines: rejected at
every normalization; its exponential steepening beyond $kx\sim2$ distorts
the subtype spacing in the opposite direction.

\emph{Route A (cooling-versus-mixing fractionation)}: eliminated twice
over. Theoretically, at the model's own radiation-confinement pressures
the cooling times at the build interface are $10^{-2}$--$10^{0}$ yr
across the plausible parameter space, while macroscopic clumps mix over
decades: the fractionation is saturated by 2--5 orders of magnitude and
cannot carry a power of $x$
\citep[cf.\ cloud-crushing and survival scalings;][]{Klein1994,GronkeOh2018}.
Empirically, Route A's unavoidable fingerprint---covering odds
proportional to the cooling function, hence to metallicity---is rejected:
multiplying the odds by $(Z/0.2\,Z_{\odot})^{\delta}$ and rerunning the
full model gives $\Delta\chi^{2}>120$ for each tested modulation
($\delta=-0.5$, $+0.5$, $+1$; Table~\ref{tab:elimination}); the law is
metallicity-blind at the tested amplitudes.

\emph{Route C (momentum-limited build against drag-limited removal)}: the
unique survivor, with its precondition verified rather than assumed. The
engine's momentum-driven lofting capacity $L_{e}/(cv_{w})$ is 14--29\%
(p16--p84) of the available nuclear mass supply, below unity for 100\% of
the selected population in every subtype: the build bottleneck is driving,
not mass, which bypasses the surface-density cancellation. The removal
channel is likewise selected by the physics: the same fast cooling that
saturates Route A makes thermal clump destruction ineffective (mixed gas
recondenses), so clumps exit the interception zone by drag entrainment,
whose timescale carries the standard density-contrast scaling
$t_{\rm drag}\propto\chi\,r/v$ with
$\chi\propto P_{\rm rad}/P_{\rm host}\propto x$. The chain
\begin{equation}
\frac{C}{1-C}=\frac{B}{G}
\propto\frac{L_{e}/(cv_{w})}{1/t_{\rm drag}}
\propto x\cdot x=x^{2}
\end{equation}
then has every structural property the confrontations demand:
scale-freeness (required by the B and D rejections), metallicity
blindness, the parity anchor, and fast relaxation. In its strongest form
the exponent is combinatorial: it counts the soft modes of the doubly
forced clump system---population (build/clear) and residence
(confinement/drag)---each equilibrating linearly in the field ratio, with
opacity frozen ($\tau\gg1$; a soft opacity mode would imprint the
rejected Route-D curvature) and stage survival not independent (its
distinctness would imprint the rejected Route-B roll-over). The
confrontations select this count among the candidates considered: the
scan disfavors $\gamma=1.5$ and $\gamma=3$ alongside the B and D
functional forms (a selection under a descriptive objective, not a
probabilistic measurement). The one imported ingredient is the
$\chi^{1}$ drag-time scaling, standard in the cloud-acceleration
literature; a targeted simulation of a radiation-confined clump in a
host-driven wind remains the external anchor.

\begin{deluxetable}{lcc}
\tablecaption{The elimination program for the covering-law exponent. All
spectral statistics are the flux-stacked line-ratio $\chi^{2}$ over eight
observables under the identical fiducial metric of
Section~\ref{sec:spectra} ($\epsilon=0.1$); joint adds the $z\le7.5$
demographic term.\label{tab:elimination}}
\tablehead{\colhead{candidate law} & \colhead{spectral $\chi^{2}$} &
\colhead{joint}}
\startdata
pure $x^{2}$ ($\gamma{=}2$, $k{=}1$) & 10.4 & 11.5\\
$\gamma=1.5$ ($k{=}1$) & 76.1 & 77.3\\
$\gamma=2.5$ ($k{=}1$) & 86.2 & 87.3\\
$\gamma=3$ ($k{=}1$) & 211.4 & 212.7\\
Route B, $k{=}1$ (parameter-free) & 165.1 & 166.6\\
Route B, best $k$ ($k{=}3$) & 27.3 & 28.9\\
Route D, best $k$ ($k{=}1.68$) & 51.9 & 53.2\\
metallicity-modulated, $\delta{=}{+}1$ & 132.9 & 135.8\\
metallicity-modulated, $\delta{=}{-}0.5$ & 241.1 & 246.7\\
\enddata
\tablecomments{The $k=1$ profile in $\gamma$ is non-monotonic
(76.1, 10.4, 86.2, 211.4 at $\gamma=1.5$, 2, 2.5, 3), with a minimum at
$\gamma=2$: changing the law reshuffles
which sources are selected into each subtype, so the objective is not a
smooth function of the exponent. The preference for $(2,1)$ over every
alternative nevertheless exceeds the seed scatter by an order of
magnitude.}
\end{deluxetable}

\section{Joint confrontations}\label{sec:results}

\emph{Statistical conventions.} We report unreduced $\chi^{2}$ sums under a
declared error model and do not convert them to likelihoods, $p$-values, or
$\Delta\chi^{2}$ probabilities; the covering-law scan objective in
particular is a descriptive sum of the spectral and demographic terms
(Section~\ref{sec:law}), and parameters are selected on coarse grids
without propagated posteriors (Section~\ref{sec:limits}). The spectral
statistic sums $N=8$ stacked observables---the [O~III]$\lambda5007$/H$\beta$
and [O~III]/[O~II] ratios of the four continuum subtypes---against the
published stack errors, with three parameters selected on this metric
($\gamma$, $k$, and the stellar ionizing coupling $\epsilon$), leaving 5
residual degrees of freedom ($\chi^{2}_{\nu}\simeq2.1$). The demographic
statistic sums the three $z\le7.5$ bins against their published errors (no
cosmic-variance term) with one refit abundance normalization, leaving 2
residual degrees of freedom ($\chi^{2}_{\nu}\simeq0.6$). The joint
objective is their unweighted sum ($N=11$ observables, 4 selected
parameters). Reduced values are indicative only.

\subsection{The stacked line ratios}\label{sec:spectra}
The decisive spectral test is performed directly in observable space. For
each subtype we form the luminosity-weighted stack of per-source line
yields (mirroring how observed stacks average spectra) and compare the
resulting [O~III]$\lambda5007$/H$\beta$ and [O~III]/[O~II] ratios with the
eight stacked measurements. One physical ingredient beyond
Equation~\ref{eq:law} is required: the fraction $\epsilon$ of surviving
stellar ionizing photons that reaches the nuclear-scale diffuse gas.
Diagnosing the raw partition shows the xLRD diffuse budget would otherwise
be dominated by stellar leak-through at a level six times too high, while
no $(\gamma,k)$ can compensate (an 18-point grid fails uniformly in ratio
space with $\epsilon=1$). Physically, most stellar Lyman-continuum photons
are consumed in local H~II regions; the canonical escape fractions of
$5$--$20$\% \citep[e.g.,][]{Flury2022} bracket the fitted
$\epsilon=0.10$ (profile flat over $0.05$--$0.15$), so we identify
$\epsilon$ with the LyC escape fraction into the wider ISM and treat its
agreement as a consistency check rather than a free success. We state
plainly that $\epsilon$ was diagnosed on, and fitted to, the same eight
ratios it now helps match; it is counted among the three selected
parameters in Table~\ref{tab:ratios}. The comparison value for the
imposed closure (12.9) is evaluated on the identical catalogue and
metric; the closure sets $f_{\rm dense}$ directly from $D$ and therefore
has no $\epsilon$ degree of freedom (its own selected parameters are the
exponent and the systematics floor of its original calibration). Because
the seed range of the physical partition (5.0--15.6) straddles 12.9,
``at least as well'' is a fiducial-seed statement.

\begin{deluxetable*}{lcccc}
\tablecaption{Direct confrontation with the stacked line ratios at
$4.5\le z\le6.5$ under the fitted law with $\epsilon=0.1$. The inverted
dense fractions (final columns) are shown as model-conditional
diagnostics only.\label{tab:ratios}}
\tablehead{\colhead{subtype} & \colhead{[O III]/H$\beta$ model} &
\colhead{observed} & \colhead{$f_{\rm dense}$ model} &
\colhead{inverted (p16--p84)}}
\startdata
xLRD & 1.20 & $1.10\pm0.10$ & 0.95 & 0.981 (0.978--0.983)\\
plusLRD & 3.00 & $3.20\pm0.10$ & 0.91 & 0.901 (0.884--0.911)\\
minusLRD & 3.98 & $4.40\pm0.30$ & 0.75 & 0.738 (0.625--0.807)\\
bLRD & 4.59 & $4.90\pm0.50$ & 0.45 & 0.086 (0.000--0.658)\\
\enddata
\tablecomments{Total $\chi^{2}=10.4$ over the eight
[O III]/H$\beta$ and [O III]/[O II] observables (seed range 5.0--15.6),
with three selected parameters $(\gamma,k,\epsilon)$; the imposed $D^{2}$
closure evaluated on the identical catalogue and metric gives 12.9. All
four He~II$\lambda4686$/H$\beta$ predictions are $\le0.015$, below the
population gate of 0.1. The bLRD inverted interval is unconstrained and
its central value is not matched; only the three red subtypes
discriminate.}
\end{deluxetable*}

Table~\ref{tab:ratios} shows the result: the raw photon partition under
the fitted covering law matches the stacked ratios at least as well as
the imposed quadratic closure evaluated on the identical metric
($\chi^{2}=10.4$ versus $12.9$), and passes the independent He~II gate.
Two disclosures accompany this. First, in the earlier inversion-space
metric the imposed closure scores $0.38$ versus the physical partition's
$1.6$; that metric's $0.05$ systematic floor---20 times the xLRD
statistical interval---controls those values, which is why we adopt the
direct ratio-space confrontation as fiducial. Second, the bLRD central
value is matched by neither model and its interval does not discriminate;
headline agreement claims rest on the three red subtypes.

\subsection{Demography}\label{sec:demog}
On the three $z\le7.5$ bins (the $z\simeq9.5$ bin is
selection-dominated; Section~\ref{sec:highz}), the demographic
$\chi^{2}$ improves from 6.90 (Thomson-depth covering) to 1.15 under the
fitted law, with one refit normalization and no covering parameter tuned
on this statistic alone. Two caveats: the $z=9.5$ pull flips sign between
configurations ($+2.3\rightarrow-2.1\sigma$), partly because the
burst-channel rate was originally tuned under the superseded covering
model (a 6--60 Gyr$^{-1}$ rate scan leaves all conclusions unchanged);
and the quoted bin errors are the published ones, without an added cosmic
variance term. The model's selected xLRD fraction at the anchor redshift
(0.2\% under the UV gate) remains a factor $\sim$60 below the observed
stack-sample share, a selection-layer failure reported in the closure
analysis and unchanged here. We also note that the demographic shape
itself is contested: a recent analysis reports no evolution in the LRD
number density from cosmic dawn to cosmic noon \citep{Loiacono2026}, in
tension with the rise--peak--decline bins fitted here, while a
spectroscopic census across $3<z<6$ \citep{Zhuang2025} and an analysis of
the $z<3$ decline \citep{Zhang2026} further underscore that the
demographic shape is observationally unsettled---a disagreement
consistent with our conclusion (Section~\ref{sec:highz}) that current
LRD demography is dominated by selection calibration.

\subsection{Broad-line widths}\label{sec:widths}
Rerunning the public RUBIES width confrontation \citep{Hviding2025} under
the fitted law improves the distribution-level agreement (weighted CDF
distance 0.354 versus 0.413; a descriptive statistic---no null
distribution is available---computed against broad-component fits that
were censored at 2500 km s$^{-1}$) and changes its interpretation
(Figure~\ref{fig:rubies}): the matched population shifts to higher
black-hole masses (median $\log M_{\rm BH}$ 6.85$\rightarrow$7.36) and
Compton-thick Thomson depths ($\tau_{\rm e}\simeq2.3$), so
$\simeq1650~{\rm km\,s^{-1}}$ of the width budget is electron-scattering
broadening---consistent with envelope models in which scattering shapes
the Balmer profiles \citep{Naidu2025,deGraaff2025}---and the descriptive
compact scale moves from $1.2\times10^{3}$ to $5.9\times10^{3}$ AU
(bootstrap 5.5--7.1$\times10^{3}$; the scan support was widened from the
published 5000 AU edge to keep this fit interior, which we disclose as a
boundary effect of the original support). The factor-of-five scale shift
between covering models quantifies a systematic on any compact-scale
inference. We note an unresolved geometric tension: the scattering
envelope's calibrated outer scale (1360 AU) lies inside the descriptive
line-emitting scale, whereas a scatterer must overlie its emitter; the
single-scale envelope is evidently oversimplified, and an extended
scattering halo ($\gtrsim6\times10^{3}$ AU) is required for
self-consistency (Section~\ref{sec:limits}). The headline scattering
quantities survive this relocation: $\tau_{\rm e}$ is computed from the
nuclear gas column, not from the 1360 AU photosphere, so the
$\simeq1650~{\rm km\,s^{-1}}$ budget is unchanged; what the relocation
does alter is the covering geometry of the scattering region, which is no
longer constrained by the continuum-reprocessor calibration---an
uncertainty we estimate at the tens-of-percent level on the width
distribution shape.

Separately, the H$\alpha$ calibration marker---the offset between the
model's selected-population median H$\alpha$ luminosity and that of a
UV- and redshift-matched public JADES control sample (not an LRD
sample)---moves from $+0.89$ to $+0.64$ dex; a factor-4 residual remains
and this quantity is a consistency indicator, not a fit.

\begin{figure*}[tp]
\centering
\includegraphics[width=0.85\textwidth]{rubies_fwhm_confrontation.pdf}
\caption{RUBIES broad Balmer width confrontation under the fitted covering
law. Left: cumulative width distributions; the dotted vertical line marks
the 2500 km s$^{-1}$ ceiling imposed in the published broad-component
fits (the model CDFs are not truncated). The dashed curve shows the
superseded 20--40 pc reservoir-scale width model for reference. Right:
descriptive compact-scale calibration; the model median matches the
observed median at $5.9\times10^{3}$ AU with a substantial
electron-scattering contribution.}
\label{fig:rubies}
\end{figure*}

\section{Predictions}\label{sec:predictions}

\begin{figure*}[tp]
\centering
\includegraphics[width=0.95\textwidth]{subtype_redshift_march.pdf}
\caption{The predicted subtype march (single fiducial realization; the
seed-averaged fractions with combined binomial-plus-seed uncertainties are
given in the text and in the companion Letter). Left: subtype fractions
versus redshift, selection-weighted (solid) and with selection removed
(dashed); the trend strengthens without selection. Right: median dense
photon fraction. The monotonic rise of the red fraction is established
across seeds for $z\gtrsim4.5$; the $z\simeq3.25$ bin (effective sample
$\simeq$19) is an extrapolation below the calibration range and is
plateau-consistent.}
\label{fig:march}
\end{figure*}

\emph{(i) The subtype march.} Because $D$ grows as reservoirs feed while
hosts assemble, the subtype mixture must evolve. Seed-averaged, with
combined binomial and seed uncertainties, the selected red fraction
(xLRD+plusLRD) is $0.04\pm0.05$, $0.09\pm0.04$, $0.23\pm0.05$,
$0.49\pm0.09$, and $0.48\pm0.14$ at $z\simeq9.5$, 7.5, 5.5, 4.5, and 3.25,
while the bLRD fraction falls from $0.61\pm0.15$ to $\lesssim0.05$
(Figure~\ref{fig:march}). Removing the selection weighting strengthens
the trend, so a transparent gate cannot manufacture it; the rise is
monotonic in all four seed realizations for $z\ge4.5$. The predicted
slope between $z\simeq7.5$ and 4.5 is $0.10$--$0.14$ per unit redshift. A
full statement of this prediction, its uncertainties, its exposure to
data-side selection evolution, and its falsification thresholds is given
in the companion Letter (Pandev 2026c), which this section summarizes.
The recent discovery of two LRDs transitioning into quasars
\citep{Fu2025} provides direct evidence for evolutionary transitions
at the phase exit, complementary to the within-phase march predicted
here.

\emph{(ii) X-ray and variability contrast.} Blue LRDs carry unbuilt
envelopes, so at matched bolometric luminosity they should be
systematically more X-ray transmitting and more variable than red LRDs,
with a steepened variability gradient across the subtype sequence. Only
the ordering is claimed; published LRD monitoring finds low variability
overall, with tentative variability in a minority of sources in larger
multi-epoch samples \citep{Tee2024,Secunda2025,KokuboHarikane2024}, and
the prediction concerns the
differential between subtypes at matched luminosity, which existing
samples have not yet tested.

\emph{(iii) Within-subtype scatter.} A deterministic $D^{2}$ law predicts
a tight line-ratio locus at fixed continuum class; the covering chain
predicts scatter from the covering lognormal. Because the scatter
amplitude is unconstrained by present data (Section~\ref{sec:law}), the
prediction is a range; object-level [O~III]/H$\beta$ dispersion within
subtypes discriminates the two pictures.

\emph{(iv) Saturation homogeneity and no hysteresis.} Covering saturates
(reaches the 0.995 cap) near $x\simeq14$, so the reddest sources should
be homogeneous in covering-driven marks; and no thick loop should appear
in the variability--color plane---the fast-equilibration hierarchy of
Section~\ref{sec:equil} forbids it.

\emph{(v) Scale-free granularity.} The soft-mode construction of
Section~\ref{sec:exponent} is self-similar: covering-driven fluctuations
(absorption variability, X-ray transmission granularity in partially
covered sources) should be scale-free; a measured characteristic clump
size would falsify the framework independently of the lifecycle.

\section{The $z\simeq9.5$ bin is a completeness thermometer}
\label{sec:highz}
The fitted model underpredicts the reported $z\simeq9.5$ broad-selection
density by a factor 6.3, a residual robust to the burst-rate scan. Its
character is resolved by a dichotomy test on the forward catalogue. The
intrinsic $M_{\rm UV}<-18.5$ nuclear-active population at $z\simeq9.5$
exceeds the reported density by a factor $\sim$34 (a model-conditional,
order-of-magnitude statement: the estimate reweights catalogued sources
by their selection probabilities, and its headroom resides in sources
with $p_{\rm sel}=10^{-6}$--$10^{-3}$ assigned by an unvalidated analytic
gate). These sources are UV-bright, compact, and embedded, and fail
selection only through the V-shape color gate, because at this epoch the
covering law leaves most envelopes unbuilt and continua blue. Rescaling
that single gate within transparent limits ($\sim$0.4 mag in the redness
threshold) moves the predicted $z=9.5$ density from a factor $\sim$6
deficit to a factor $\sim$3 excess of the observed value while moving the
$z\le7.5$ bins by well under their errors. Three conclusions follow. No
additional formation channel is required by current number densities. The
$z\simeq9.5$ bin cannot presently discriminate formation physics; it
measures the survey's color completeness, which \citet{Rinaldi2026}
describe as essentially unconstrained at $z>6$--7 (their discussion of
MIRI depth). And the falsifiable content is compositional: the real
$z>8.5$ LRD population should be dominated by blue-continuum,
low-covering compact sources (quantified in the Letter, including the
per-subtype gate-passing fractions that reconcile bLRD dominance of the
selected sample with the removal of the blue headroom); a red-dominated
census would falsify this resolution. The decisive external products are
injection/recovery completeness for blue compact templates and the
effective-area masks.

\section{Limitations}\label{sec:limits}
The observation layer remains a set of transparent proxies: the selection
gates are not the survey pipelines; the X-ray transfer over-suppresses
known bands; and the Balmer decrement is not solved (intrinsic
H$\alpha$/H$\beta\simeq2.8$--3.2 against observed 5.5--10.1
\citep{PerezGonzalez2026}; a smooth screen with $A_{V}\simeq2$--4 toward
the line region is not rejected). The covering-law parameters are
selected on coarse grids without propagated posteriors; the covering
scatter is assumed; $\epsilon$ is fitted (its agreement with canonical
escape fractions is a consistency, not a derivation); and the drag-time
scaling of Section~\ref{sec:exponent} is imported from standard cloud
physics. The single-scale scattering envelope is internally inconsistent
with the descriptive line-emitting scale and requires an extended halo.
The published RUBIES widths are censored at $2500~{\rm km\,s^{-1}}$, so
the compact-scale calibration is descriptive. Demographic errors omit
cosmic variance. Seed statements rest on four realizations. None of the
fits here is a posterior; a catalogue-level likelihood with
injection/recovery completeness remains the requirement for population
inference.

\section{Summary}\label{sec:summary}
A single equilibrium relation---covering odds equal to the squared
intrinsic engine-to-host contrast, with a momentum-parity
normalization---simultaneously (1) matches the stacked line-ratio
sequence through a raw photon partition at least as well as an imposed
closure, given one geometric parameter consistent with canonical LyC
escape fractions, (2) improves the reliable-bin demographic fit to
$\chi^{2}=1.2$ on 3 bins, (3) improves and reinterprets the public
broad-width confrontation in favor of electron-scattering broadening, and
(4) converts the high-redshift abundance residual into a quantified
statement about color completeness. The exponent is constructed, not
fitted: of four candidate mechanisms confronted, three were eliminated by
the data (with the eliminations consistent with the clump system having
exactly two soft modes), and the surviving candidate's
precondition---momentum-limited building---holds for the entire selected
population. The framework's failure modes---a red-dominated $z>8.5$
census, a tight within-subtype line-ratio locus, an inverted subtype
march, a thick variability--color hysteresis loop, or a measured
characteristic clump scale---are all measurable with data expected on a
1--2 year horizon.

\begin{acknowledgments}
This work is unfunded and was carried out independently. It made use of
public data products from the JWST archive and of published catalogues by
the teams cited herein; the author thanks those teams for making their
data and measurements public. This research made use of NASA's
Astrophysics Data System and the arXiv preprint repository.
\end{acknowledgments}

\facilities{JWST (NIRSpec, NIRCam, MIRI), Chandra (data limits as cited)}

\section*{Data and code availability}
The lifecycle model, covering-law scans, Cloudy anchor grids, stack
inversion, elimination-program scripts, and all confrontation scripts
accompany this paper as an MIT-licensed code archive, together with
machine-readable results for every table and figure (Zenodo DOI to be
assigned at submission); the full 65-configuration forward-catalogue grid
is deposited as a separate CC-BY-4.0 data record (Zenodo DOI to be
assigned). Only
public data products are used: the \citet{Rinaldi2026} and \citet{Ma2026}
binned densities, the \citet{PerezGonzalez2026} stacks, the
\citet{Hviding2025} RUBIES release, the \citet{Sok2026} hardness
constraint, and public JADES catalogues \citep{Eisenstein2023}.

\software{Cloudy \citep{Chatzikos2023}, numpy, scipy, matplotlib,
astropy, colossus}

\begin{thebibliography}{}
\bibitem[Akins et al.(2024)]{Akins2024} Akins, H. B., et al. 2024,
arXiv:2406.10341
\bibitem[Ananna et al.(2024)]{Ananna2024} Ananna, T. T., et al. 2024,
ApJL, 969, L18
\bibitem[Barro et al.(2025)]{Barro2025} Barro, G., et al. 2025,
arXiv:2512.15853
\bibitem[Begelman \& Dexter(2026)]{Begelman2026} Begelman, M. C., \&
Dexter, J. 2026, ApJ, 996, 48
\bibitem[Bullock et al.(2001)]{Bullock2001} Bullock, J. S., Dekel, A.,
Kolatt, T. S., et al. 2001, ApJ, 555, 240
\bibitem[Carranza-Escudero et al.(2025)]{Carranza2025} Carranza-Escudero,
M., et al. 2025, ApJL, 989, L50
\bibitem[Chatzikos et al.(2023)]{Chatzikos2023} Chatzikos, M., et al.
2023, RMxAA, 59, 271
\bibitem[de Graaff et al.(2025)]{deGraaff2025} de Graaff, A., et al.
2025, arXiv:2503.16600
\bibitem[de Graaff et al.(2025b)]{deGraaff2025b} de Graaff, A., et al.
2025, arXiv:2511.21820
\bibitem[Eisenstein et al.(2023)]{Eisenstein2023} Eisenstein, D. J., et
al. 2023, arXiv:2306.02465
\bibitem[Fakhouri et al.(2010)]{Fakhouri2010} Fakhouri, O., Ma, C.-P., \&
Boylan-Kolchin, M. 2010, MNRAS, 406, 2267
\bibitem[Flury et al.(2022)]{Flury2022} Flury, S. R., et al. 2022, ApJS,
260, 1
\bibitem[Fu et al.(2025)]{Fu2025} Fu, S., et al. 2025, arXiv:2512.02096
\bibitem[Greene et al.(2024)]{Greene2024} Greene, J. E., et al. 2024,
ApJ, 964, 39
\bibitem[Gronke \& Oh(2018)]{GronkeOh2018} Gronke, M., \& Oh, S. P. 2018,
MNRAS, 480, L111
\bibitem[Hviding et al.(2025)]{Hviding2025} Hviding, R. E., et al. 2025,
RUBIES broad-line data release, Zenodo, doi:10.5281/zenodo.16758549
\bibitem[Kido et al.(2025)]{Kido2025} Kido, D., et al. 2025, MNRAS, 544,
3407
\bibitem[Kim \& Ostriker(2015)]{KimOstriker2015} Kim, C.-G., \& Ostriker,
E. C. 2015, ApJ, 802, 99
\bibitem[Klein et al.(1994)]{Klein1994} Klein, R. I., McKee, C. F., \&
Colella, P. 1994, ApJ, 420, 213
\bibitem[Kocevski et al.(2024)]{Kocevski2024} Kocevski, D. D., et al.
2024, arXiv:2404.03576
\bibitem[Kocevski et al.(2025)]{Kocevski2025} Kocevski, D. D., et al.
2025, ApJ, 986, 126
\bibitem[Inayoshi et al.(2025)]{Inayoshi2024} Inayoshi, K., Kimura, S. S.,
\& Noda, H. 2025, PASJ, 77, 811
\bibitem[Kokubo \& Harikane(2025)]{KokuboHarikane2024} Kokubo, M., \&
Harikane, Y. 2025, ApJ, 995, 24
\bibitem[Loiacono et al.(2026)]{Loiacono2026} Loiacono, F., et al. 2026,
arXiv:2606.30253
\bibitem[Ma et al.(2026)]{Ma2026} Ma, Y., et al. 2026, ApJ, 1000, 59
\bibitem[Maiolino et al.(2024)]{Maiolino2024} Maiolino, R., et al. 2024,
arXiv:2405.00504
\bibitem[Matthee et al.(2024)]{Matthee2024} Matthee, J., et al. 2024,
ApJ, 963, 129
\bibitem[Mo et al.(1998)]{Mo1998} Mo, H. J., Mao, S., \& White, S. D. M.
1998, MNRAS, 295, 319
\bibitem[Naidu et al.(2025)]{Naidu2025} Naidu, R. P., et al. 2025,
arXiv:2503.16596
\bibitem[Pacucci \& Loeb(2025)]{PacucciLoeb2025} Pacucci, F., \& Loeb,
A. 2025, ApJL, arXiv:2506.03244
\bibitem[Pacucci \& Narayan(2024)]{PacucciNarayan2024} Pacucci, F., \&
Narayan, R. 2024, ApJ, arXiv:2407.15915 %% verify arXiv ID on ADS
\bibitem[P\'erez-Gonz\'alez et al.(2026)]{PerezGonzalez2026}
P\'erez-Gonz\'alez, P. G., et al. 2026, arXiv:2602.20247
\bibitem[Pizzati et al.(2025)]{Pizzati2025} Pizzati, E., et al. 2025,
MNRAS, 539, 2910
\bibitem[Rinaldi et al.(2026)]{Rinaldi2026} Rinaldi, P., et al. 2026,
arXiv:2604.07138
\bibitem[Rusakov et al.(2025)]{Rusakov2025} Rusakov, V., et al. 2025,
Nature, arXiv:2503.16595 %% verify author list on ADS
\bibitem[Sacchi \& Bogd\'an(2025)]{Sacchi2025} Sacchi, A., \& Bogd\'an,
\'A. 2025, ApJL, 989, L30
\bibitem[Santarelli et al.(2025)]{Santarelli2025} Santarelli, A. D., et
al. 2025, arXiv:2510.17952
\bibitem[Secunda et al.(2025)]{Secunda2025} Secunda, A., et al. 2025,
arXiv:2509.03571
\bibitem[Setton et al.(2024)]{Setton2024} Setton, D. J., et al. 2024,
arXiv:2411.03424
\bibitem[Sneppen et al.(2026)]{Sneppen2026} Sneppen, A., et al. 2026,
arXiv:2605.15263
\bibitem[Sneppen et al.(2026b)]{Sneppen2026b} Sneppen, A., et al. 2026,
arXiv:2601.18864
\bibitem[Sok et al.(2026)]{Sok2026} Sok, V., et al. 2026, arXiv:2606.23778
\bibitem[Sun et al.(2026)]{Sun2026} Sun, W. Q., et al. 2026, The Open
Journal of Astrophysics, 9
\bibitem[Tee et al.(2024)]{Tee2024} Tee, W. L., et al. 2024,
arXiv:2412.05242
\bibitem[Tinker et al.(2008)]{Tinker2008} Tinker, J., et al. 2008, ApJ,
688, 709
\bibitem[Tinker et al.(2010)]{Tinker2010} Tinker, J., et al. 2010, ApJ,
724, 878
\bibitem[Tortosa et al.(2026)]{Tortosa2026} Tortosa, A., et al. 2026,
A\&A, 708, A293
\bibitem[Wang et al.(2026)]{Wang2026} Wang, Z., et al. 2026,
arXiv:2603.15736
\bibitem[Yue et al.(2024)]{Yue2024} Yue, M., et al. 2024, arXiv:2404.13290
\bibitem[Zhang et al.(2026)]{Zhang2026} Zhang, C., et al. 2026,
arXiv:2606.02773
\bibitem[Zhuang et al.(2025)]{Zhuang2025} Zhuang, M.-Y., et al. 2025,
arXiv:2505.20393
\end{thebibliography}

\end{document}
